\noindent
\includegraphics[height=1.25cm]{images/pictograms/replication}
\includegraphics[height=1.25cm]{images/pictograms/aspect_logo}
\includegraphics[height=1.25cm]{images/pictograms/ice}
\includegraphics[height=1.25cm]{images/pictograms/visualisation}
\includegraphics[height=1.25cm]{images/pictograms/gravity}
\includegraphics[height=1.25cm]{images/pictograms/elasticity}
\includegraphics[height=1.25cm]{images/pictograms/benchmark}
\includegraphics[height=1.25cm]{images/pictograms/triangle}
\includegraphics[height=1.25cm]{images/pictograms/under_construction}
\includegraphics[height=1.25cm]{images/pictograms/bsc}
\includegraphics[height=1.25cm]{images/pictograms/msc}
\includegraphics[height=1.25cm]{images/pictograms/tools}
\includegraphics[height=1.25cm]{images/pictograms/FEM}
\includegraphics[height=1.25cm]{images/pictograms/FDM}
\includegraphics[height=1.25cm]{images/pictograms/temperature}
\includegraphics[height=1.25cm]{images/pictograms/3d}
\includegraphics[height=1.25cm]{images/pictograms/pic}
\includegraphics[height=1.25cm]{images/pictograms/nonlinear}
\includegraphics[height=1.25cm]{images/pictograms/paraview}
\includegraphics[height=1.25cm]{images/pictograms/publication}
\includegraphics[height=1.25cm]{images/pictograms/streamfunction}
\includegraphics[height=1.25cm]{images/pictograms/wave}

%%%%%%%%%%%%%%%%%%%%%%%%%%%%%%%%%%%%%%%%%%%%%%%%%%%%%%%%%%%%%%%%%%%%%%%%%%%%%%%%%%%%%%%%%%%%%%%%%%%

\begin{flushright} {\tiny {\color{gray} python\_codes/fieldstone\_170/text.tex}} \end{flushright}

%\lstinputlisting[language=bash,basicstyle=\small]{python_codes/template_keywords.key}

\par\noindent\rule{\textwidth}{0.4pt}

\begin{center}
\inpython
{\small Code: \url{https://github.com/cedrict/fieldstone/tree/master/python_codes/fieldstone_170}}
\end{center}

\par\noindent\rule{\textwidth}{0.4pt}

{\sl The code for this stone was developed by Djurre van der Molen}. \index{contributors}{D. van der Molen}

\par\noindent\rule{\textwidth}{0.4pt}

Last revision: March 6th, 2025.

\par\noindent\rule{\textwidth}{0.4pt}

%%%%%%%%%%%%%%%%%%%%%%%%%%%%%%%%%%%%%%%%%%%%%%%%%%%%%%%%%%%%%%%%%%%%%%%%%%%%%%%%%%%%%%%%%%%%%%%%%%%

The idea behind this \stone originates in \stone~126 in which I attempt to replicate a 
published article from 2003 \cite{grfr03}. Unfortunately, upon my asking, 
the authors were able to confirm that the original code and results were lost for good. 
This means that my replication hinges on how my my figures match 
the figures of the paper such as the following one:

\begin{center}
\includegraphics[width=10cm]{python_codes/fieldstone_170/images/grfr03}
\end{center}

This is quite problematic because a) the quality of the image is objectively not great; 
b) the colorscale used here is not standard (it is not available in ParaView for example); 
c) there could have been color distortion from the original during the publication 
process and the later digitization in order to publish the article online as a pdf document.

Let us then isolate one panel and the color scale:

\begin{center}
\includegraphics[height=4cm]{python_codes/fieldstone_170/images/graph_3_full}
\includegraphics[height=4cm]{python_codes/fieldstone_170/images/legenda}
\end{center}

In order to carry out a more quantitative comparison
we first need to read in (digitize) the left panel and transform each pixel into 
a number (or several numbers such as red, green, blue values), and second use the 
color bar to find the corresponding pressure value between 0 and 1.

The code showcases five functions:

\begin{itemize}
\item \lstinline{generate_color_value_map(image_path)}:
\item \lstinline{get_value_from_color(color,color_value_map,color_tree)}:
\item \lstinline{pixel_to_data_coords(x_pixel,y_pixel,img_width,img_height,x_range,y_range)}:
\item \lstinline{analyze_graph_image(image_path,color_value_map,color_tree,x_range,y_range)}:
\item \lstinline{plot_3d_surface(x_vals,y_vals,z_vals)}:
\end{itemize}




pip install opencv-python

